\begin{enumerate}
    \item[\textbf{Problema 5.1.}] Se você fosse solicitado a conduzir um estudo populacional, o que você estaria procurando, o que seria conhecido e o que você gostaria de saber?

    \item[\textbf{Problema 5.2.}] Que tipo de suposições você faria se fosse solicitado a realizar um estudo populacional? Com base em que essas suposições podem ser justificadas?

    \item[\textbf{Problema 5.3.}] Que fatores podem restringir ou influenciar o crescimento desenfreado visto nas Figuras 5.1 e 5.2?

    \item[\textbf{Problema 5.4.}] Determine o raio da Terra em metros e pés a partir das áreas de superfície indicadas na Seção 5.1. Esses valores são consistentes com os fatores de conversão fornecidos na Tabela 2.3?

    \item[\textbf{Problema 5.5.}] Confirme que levaria quase duzentos mil anos para contar uma população de 5,63 milhões de bilhões de pessoas a uma taxa de 1000 pessoas/s.

    \item[\textbf{Problema 5.6.}] Confirme, derivando a eq. (5.4), que o $N(t)$ nela dado satisfaz a eq. (5.1).

    \item[\textbf{Problema 5.7.}] Como as projeções das Figuras 5.1 e 5.2 mudariam se a taxa de crescimento fosse, respectivamente, de 1\% ao ano e 3\% ao ano?

    \item[\textbf{Problema 5.8.}] Que taxa de crescimento anual seria necessária para dobrar o dinheiro em sete anos?

    \item[\textbf{Problema 5.9.}] Verifique se a eq. (5.12) é a solução para a equação diferencial exponencial conforme dada na eq. (5.10) usando o resultado de que $\int du/u = \ln u + \text{constante}$.

    \item[\textbf{Problema 5.10.}] Mostre que a solução (5.12) para a equação diferencial (5.10) também pode ser encontrada assumindo a seguinte solução experimental para $N(t)$: $N(t) = Ce^{\alpha t}$.

    \item[\textbf{Problema 5.11.}] Por que a constante de proporcionalidade na eq. (5.15) é equivalente a ter $\lambda < 0$ na eq. (5.10)? Que tipo de comportamento esperaríamos então?

    \item[\textbf{Problema 5.12.}] Verifique se cada termo na eq. (5.22) tem as mesmas dimensões físicas.

    \item[\textbf{Problema 5.13.}] Confirme se a eq. (5.23) é a solução correta para a eq. (5.22).

    \item[\textbf{Problema 5.14.}] Use as definições de resistência e capacitância para verificar se o produto $RC$ tem as dimensões físicas do tempo.

    \item[\textbf{Problema 5.15.}] Verifique se a eq. (5.29) foi derivada corretamente da eq. (5.28).

    \item[\textbf{Problema 5.16.}] Aplique a Lei das Malhas de Kirchhoff (KVL) da eq. (5.28) ao circuito na Figura 5.7 para confirmar a validade da eq. (5.27).

    \item[\textbf{Problema 5.17.}] Confirme, derivando a eq. (5.31), que ela satisfaz a eq. (5.29).

    \item[\textbf{Problema 5.18.}] Calcule a inclinação inicial do carregamento do capacitor a partir da solução indicada na eq. (5.31).

    \item[\textbf{Problema 5.19.}] As eq. (5.35) e (5.38) estão relacionadas? Como?

    \item[\textbf{Problema 5.20.}] Verifique todas as etapas que levam da eq. (5.35) à eq. (5.39).

    \item[\textbf{Problema 5.21.}] Construa uma versão da Tabela 5.3 para uma taxa de juros anual de 18\%.

    \item[\textbf{Problema 5.22.}] Se a gasolina custava US$ 0,70/galão em 1978, calcule e compare as taxas de inflação de preços para os intervalos 1933–1978 e 1978–2001.

    \item[\textbf{Problema 5.23.}] Se o custo do dinheiro exceder o custo das mercadorias, o que acontece com $\lambda_{real}$?

    \item[\textbf{Problema 5.24.}] Especule sobre os efeitos potenciais de $\lambda_{real}$ permanecer negativo por longos períodos de tempo.

    \item[\textbf{Problema 5.25.}] Pesquise os números de nascimento, morte, imigração e emigração dos EUA para as 10 décadas do século 20 e use a equação de equilíbrio (5.45) para calcular as mudanças populacionais que essas taxas preveem.

    \item[\textbf{Problema 5.26.}] Como as previsões do Problema 5.25 se comparam com os dados reais do censo dos EUA?

    \item[\textbf{Problema 5.27.}] Quais são as implicações para o modelo da eq. (5.48) ao afrouxar a restrição de que $\lambda_1, \lambda_2 > 0$?

    \item[\textbf{Problema 5.28.}] Confirme, derivando a eq. (5.52), que ela satisfaz a eq. (5.51).

    \item[\textbf{Problema 5.29.}] Assumindo que a eq. (5.54b) pode ser resolvida para $F(t)$, mostre que $E(t)$ pode ser determinado sem integração adicional.

    \item[\textbf{Problema 5.30.}] Confirme se a eq. (5.57) segue da eq. (5.55) via eq. (5.56).

    \item[\textbf{Problema 5.31.}] Suponha que você receba a solução para a população inimiga que satisfaça a eq. (5.53) como $E(t) = E_0 \cosh \alpha t - \sqrt{\lambda_F / \lambda_E} F_0 \sinh \alpha t$, onde $\alpha^2 = \lambda_E \lambda_F$. Quanto tempo leva para as forças inimigas serem completamente aniquiladas?

    \item[\textbf{Problema 5.32.}] A estratégia de dividir e conquistar funcionaria para uma "lei de atrito linear" que, para exércitos igualmente eficazes, substitui a eq. (5.57) por $F(t) - E(t) = F_0 - E_0$?

    \item[\textbf{Problema 5.33.}] Mostre que se levar um tempo, $t_c$, para contar uma população, $P(t)$, que tem uma taxa de crescimento de $\lambda$, a população aumentará em uma quantidade igual a $\lambda t_c P(t)$.

    \item[\textbf{Problema 5.34.}] (a) Se a taxa de contagem da população for $c$, quanto tempo leva para contar a população no tempo $t$? (b) Quanto tempo leva para contar o aumento da população que ocorreu enquanto ela estava sendo contada no tempo $t$?

    \item[\textbf{Problema 5.35.}] Encontre os números reais da população mundial para 1970, 1980, 1990 e 2000. Use esses dados para atualizar as projeções mostradas nas Figuras 5.1 e 5.2.

    \item[\textbf{Problema 5.36.}] Usando a taxa de crescimento de 2\%, plote a população mundial de 1960 a 2060 com uma escala de ordenadas de 10 bilhões de pessoas por 1,50 polegadas. Essa curva parece um crescimento exponencial "razoável"?

    \item[\textbf{Problema 5.37.}] As escalas de ordenadas das Figuras 5.1 e 5.2 são, respectivamente, 1,5 pol = 100 bilhões de pessoas e 1,5 pol = 3 milhões de bilhões de pessoas. Quanto papel é necessário para traçar a população mundial de 1960 (3 bilhões de pessoas) e a população mundial projetada de 2692 (5,63 × $10^{15}$ pessoas) usando cada uma dessas escalas?

    \item[\textbf{Problema 5.38.}] Trace o crescimento da população mundial de um valor inicial em 1960 de 3 bilhões de pessoas a taxas de crescimento de 1\%, 2\% e 3\% ao ano até o ano 2700 usando papel semi-logarítmico. Que formas têm essas curvas? Quais são suas inclinações e interceptações?

    \item[\textbf{Problema 5.39.}] Qual é a constante de tempo, comparável à de um circuito RC, para uma população que decai a uma taxa fracionária por unidade de tempo $\lambda$?

    \item[\textbf{Problema 5.40.}] Quanto deve ser reservado em 2002 numa conta poupança a 5,5\% ao ano para acumular US$ 1.000.000 até 2022? Até 2042?

    \item[\textbf{Problema 5.41.}] Suponha que houvesse uma taxa de inflação constante de 3\% ao ano, quais deveriam ser os investimentos do Problema 5.40 para acumular US$ 1.000.000 em 2042, medidos em dólares de 2002?

    \item[\textbf{Problema 5.42.}] O notável historiador Stephen Ambrose narrou o crescimento das ferrovias americanas listando as seguintes quantidades de total de trilhos por década: 726 mi (1834), 4.311 (1844), 15.675 (1854) e 33.860 (1864). Determine: (a) a taxa de crescimento década a década; e (b) a taxa de crescimento exponencial considerando todos os dados fornecidos.

    \item[\textbf{Problema 5.43.}] Verifique, por diferenciação e substituição, se a solução a seguir satisfaz as eqs. (5.53): $F(t) = F_0 \cosh \alpha t - \sqrt{\lambda_E / \lambda_F} E_0 \sinh \alpha t$ e $E(t) = E_0 \cosh \alpha t - \sqrt{\lambda_F / \lambda_E} F_0 \sinh \alpha t$.

    \item[\textbf{Problema 5.44.}] Confirme se a solução verificada no Problema 5.43 satisfaz a lei do quadrado de Lanchester da eq. (5.58).

    \item[\textbf{Problema 5.45.}] As forças iniciais de dois exércitos adversários são $F_0 = 10.000$ e $E_0 = 5.000$ soldados, com taxas de perda iguais de 0,1 por dia. Quem vai ganhar? Quanto tempo vai demorar a batalha? (Dica: consulte o problema 5.31.) Quantas tropas terá o vencedor quando o inimigo for derrotado? Faça um gráfico das populações do exército até que o inimigo seja completamente aniquilado.

    \item[\textbf{Problema 5.46.}] As forças iniciais de dois exércitos adversários são $F_0 = 10.000$ e $E_0 = 5.000$ soldados, e $\lambda_F = 0,1$ por dia. Quem vai ganhar e com quais forças restantes se $\lambda_E = 0,2$, $0,5$ e $1,0$ por dia? Que valor de $\lambda_E$ produziria um empate?

    \item[\textbf{Problema 5.47.}] A batalha histórica da Segunda Guerra Mundial em Iwo Jima começou com tropas de $F_0 = 54.000$ e $E_0 = 21.500$ soldados, com $\lambda_F = 0,0106$ por dia e $\lambda_E = 0,0544$ por dia. Na ausência de reforços, quanto tempo essa batalha teria durado? Quantas tropas o vencedor teria quando as forças do perdedor estivessem totalmente esgotadas?

    \item[\textbf{Problema 5.48.}] Para acabar com a luta por Iwo Jima em 28 dias, quantas tropas os Estados Unidos teriam que ter inicialmente? Como as perdas dos EUA neste cenário se comparam às encontradas no cenário do Problema 5.47?
\end{enumerate}
